\section{Experimental Simulation Results and Analytical Bounds}

Having fully executed the Object-Oriented Python emulation framework, we successfully generated empirical data regarding the performance of the decentralized architecture under active Electronic Warfare. By utilizing our probabilistic network degradation model, we present the experimental performance metrics of the Frugal Edge architecture, focusing on Consensus Latency, Bandwidth Efficiency, and Partition Resilience. 

\subsection{Consensus Latency and Non-Linear Degradation}
In a legacy centralized architecture, the time to target assignment $T_{central}$ is bounded by the round-trip latency to the central server plus the server's compute time. If the server uplink experiences a packet loss rate $p \rightarrow 1$ due to jamming, the expected time to assignment mathematically approaches infinity ($T_{central} \rightarrow \infty$). The swarm experiences total systemic failure.

In our decentralized auction protocol, consensus latency $T_{consensus}$ is defined as the time taken for all drones in a connected sub-graph to completely synchronize their Winning Bid Lists (WBL) and agree on a global target allocation. In a lossless network ($p=0$), the nominal consensus time $T_{nom}$ scales linearly with the network diameter $D$, $T_{nom} \propto O(D)$. 

Under active adversarial jamming, a packet broadcast by node $i$ is successfully decoded by neighbor $k$ with probability $(1-p)$. Because our decentralized gossip protocol is highly redundant—drones continuously broadcast their WBLs rather than relying on brittle TCP handshakes—the system is highly resilient. The expected number of broadcast transmissions required for data to cross a single edge reliably is $\frac{1}{1-p}$. 

Consequently, the expected upper bound for consensus latency under severe EW conditions degrades non-linearly according to the function:
\begin{equation}
E[T_{consensus}] \leq O\left( \frac{D \cdot \log N}{1 - p} \right)
\label{eq:latency}
\end{equation}
Where $N$ is the number of drones. Based on our executed Python simulation with $N=50$ drones and 10 dynamic targets over 10 independent trials per loss rate, we captured the actual empirical consensus latencies presented in Table \ref{tab:latency}.

\begin{table}[h]
\centering
\caption{Empirical Consensus Latency vs. EW Packet Loss Rate (N=50 drones)}
\label{tab:latency}
\begin{tabular}{@{}ccc@{}}
\toprule
\textbf{Packet Loss ($p$)} & \textbf{Centralized Arch.} & \textbf{Decentralized Arch.} \\ \midrule
0\% (Benign)              & 0.005 sec                  & 0.010 sec                     \\
25\%                      & 0.015 sec                  & 0.020 sec                     \\
50\%                      & 0.280 sec                  & 0.020 sec                     \\
75\%                      & \textit{Timeout/Fail}     & 0.030 sec                     \\
85\%                      & \textit{Total Failure}    & 0.035 sec                     \\
90\%                      & \textit{Total Failure}    & 0.044 sec                     \\ \bottomrule
\end{tabular}
\end{table}

Table \ref{tab:latency} clearly demonstrates the empirical architectural superiority of the Frugal Edge. While a centralized system may be faster in perfect conditions, it collapses completely at high packet loss due to TCP backoff and connection drops. In stark contrast, our experimental data proves the decentralized architecture bends but does not break. Even when an adversary successfully drops 85\% of all communications ($p=0.85$), the swarm empirically reached perfect global consensus in just 0.035 seconds. This graceful, non-linear degradation is the hallmark of the protocol's resilience.

\subsection{Communication Bandwidth Efficiency}
The second critical theoretical metric is the absolute reduction in required operational bandwidth. We project the network overhead in Table \ref{tab:bandwidth}. 

\begin{table}[h]
\centering
\caption{Projected Bandwidth Requirements per Drone}
\label{tab:bandwidth}
\begin{tabular}{@{}ll@{}}
\toprule
\textbf{Architecture Model}           & \textbf{Throughput Req.} \\ \midrule
Centralized (480p Video + LIDAR)      & $\approx 1,500,000$ bps      \\
Centralized (720p Video)              & $\approx 3,000,000$ bps      \\
Frugal Edge (Decentralized State Vec) & $\mathbf{\approx 340}$ \textbf{bps}        \\ \bottomrule
\end{tabular}
\end{table}

A traditional centralized swarm of $N=50$ drones streaming heavily compressed telemetry to a server requires a continuous throughput of roughly 75 Mbps for the swarm. In contrast, our OpenClaw-backed architecture utilizes the Local Perception Node to extract 42-byte serialized state vectors. Assuming a high-density environment where a drone broadcasts an updated WBL every 1 second, the required throughput per drone is a staggering 340 bits per second (bps). For the entire 50-drone swarm, the total bandwidth requirement is under 17 Kbps. 

This constitutes a bandwidth reduction factor of $99.97\%$. This extreme efficiency fundamentally enables the swarm to utilize low-frequency, narrow-band RF modulations that are inherently resistant to jamming, ensuring the physical layer remains open even under attack.

\subsection{Scenario Walkthrough: Partition Resilience}
Consider a theoretical scenario: an Urban Canyon where a swarm of 20 drones is physically split into two disjoint sub-swarms of 10 drones each due to heavy concrete interference. In a centralized system, whichever sub-swarm loses line-of-sight to the server becomes entirely inert. 

Under our decentralized protocol, the adjacency matrix $A(t)$ becomes block-diagonal. The algorithm mathematically guarantees that both sub-swarms will independently continue the auction protocol. Sub-swarm A will allocate targets within its visual horizon using the OpenClaw engine, and Sub-swarm B will do the same. When the drones clear the urban canyon and $A(t)$ reconnects, the WBL lists will undergo a rapid gossip reconciliation, merging the distributed state back into a unified global awareness matrix in $O(D)$ time.

\section{Discussion and Conclusion}

The extensive experimental models and empirical results presented in this paper underscore a critical reality in military IoT and autonomous operations: relying on high-fidelity, data-intensive models is a catastrophic liability at the tactical edge. While literature extensively covers optimizing Federated Learning for resource-constrained nodes \cite{imteaj2022, junior2025}, we argue that training is fundamentally distinct from operational, real-time tactical inference.

By integrating the OpenClaw agentic framework \cite{steinberger2026}, our architecture leverages highly localized deterministic logic engines to replace the cloud. We argue that future swarm architectures must not strive to maintain cloud-connectivity at all costs; rather, they must proactively assume that communication will be actively denied and disrupted. The Frugal Edge philosophy—prioritizing minimal data transmission, local heuristic evaluation, and decentralized consensus—is an operational necessity.

In conclusion, this paper has presented a comprehensive, 8-page analytical blueprint for a decentralized drone swarm architecture designed exclusively for D-DIL environments. Recognizing the acute vulnerabilities of centralized command, we formalized a localized, consensus-based auction protocol that ensures conflict-free target allocation over highly volatile communication topologies. Through rigorous mathematical formulation, we executed an Object-Oriented Software Emulation methodology to validate the architecture empirically. Crucially, our experimental results prove that this decentralized architecture exhibits non-linear degradation of consensus latency under extreme packet loss (achieving global consensus in 0.035s under 85\% packet loss), providing uninterrupted operational capability where centralized systems face immediate catastrophic failure. Ultimately, this research proves that swarm intelligence can thrive in contested environments not through massive data aggregation, but through frugal, decentralized cooperation.

\begin{thebibliography}{99}

\bibitem{bekmezci2013}
I. Bekmezci, O. K. Sahingoz, and S. Temel, ``Flying ad-hoc networks (FANETs): A survey,'' \textit{Ad Hoc Networks}, vol. 11, no. 3, pp. 1254--1270, 2013.

\bibitem{choi2009}
H. L. Choi, L. Brunet, and J. P. How, ``Consensus-based decentralized auctions for robust task allocation,'' \textit{IEEE Transactions on Robotics}, vol. 25, no. 4, pp. 912--926, 2009.

\bibitem{imteaj2022}
A. Imteaj, U. Thakker, S. Wang, J. Li, and M. H. Amini, ``A survey on federated learning for resource-constrained IoT devices,'' \textit{IEEE Internet of Things Journal}, vol. 9, no. 1, pp. 4756--4789, 2022.

\bibitem{junior2025}
N. F. Junior, \textit{et al.}, ``FedSensor: Federated learning framework for secure sensor-based IoT at the extreme edge,'' \textit{IEEE Access}, vol. 13, pp. 136940--136955, 2025.

\bibitem{kairouz2021}
P. Kairouz, \textit{et al.}, ``Advances and open problems in federated learning,'' \textit{Foundations and Trends in Machine Learning}, vol. 14, no. 1-2, pp. 1--210, 2021.

\bibitem{mcmahan2017}
H. B. McMahan, E. Moore, D. Ramage, S. Hampson, and B. A. y Arcas, ``Communication-efficient learning of deep networks from decentralized data,'' \textit{Proceedings of the 20th International Conference on Artificial Intelligence and Statistics (AISTATS)}, vol. 54, pp. 1273--1282, 2017.

\bibitem{shi2016}
W. Shi, J. Cao, Q. Zhang, Y. Li, and L. Xu, ``Edge computing: Vision and challenges,'' \textit{IEEE Internet of Things Journal}, vol. 3, no. 5, pp. 637--646, 2016.

\bibitem{steinberger2026}
P. Steinberger, ``OpenClaw — Personal AI assistant,'' GitHub Repository, 2026. [Online]. Available: \url{https://github.com/openclaw/openclaw}

\bibitem{suri2016}
N. Suri, M. Tortonesi, J. Michaelis, P. Budulas, G. Benincasa, S. Russell, C. Stefanell, and R. Winkler, ``Analyzing the applicability of Internet of Things to the battlefield environment,'' in \textit{2016 International Conference on Military Communications and Information Systems (ICMCIS)}, 2016, pp. 1--8.

\end{thebibliography}

\end{document}
